% Part: incompleteness
% Chapter: introduction
% Section: overview

\documentclass[../../../include/open-logic-section]{subfiles}

\begin{document}

\olfileid[fr]{inc}{int}{ovr}

\olsection{Vue d'ensemble des résultats d'incomplétude}

Hilbert espérait formaliser les mathématiques dans une théorie
!!{axiomatizable} dont on pourrait démontrer qu'elle est complète
et !!{decidable}. Il voulait en outre établir sa cohérence par des
moyens très faibles, dits « finitistes », afin de défendre les
mathématiques classiques contre les objections intuitionnistes.
Les théorèmes d'incomplétude de Gödel ont montré que ces objectifs
ne pouvaient être atteints ensemble.

Le premier théorème d'incomplétude de Gödel montrait qu'une version
des \emph{Principia Mathematica} de Russell et Whitehead n'est pas
complète. Mais sa preuve était assez générale pour s'appliquer à
une grande variété de théories. Le problème ne tenait donc pas
seulement à ce que les \emph{Principia Mathematica} n'avaient pas
réussi à saisir intégralement les mathématiques : aucune théorie
répondant aux conditions du théorème ne peut être complète. Il a
fallu du temps pour isoler les propriétés des théories qui suffisent
à l'application des théorèmes d'incomplétude et pour reformuler la
preuve de Gödel en fonction de ces seules propriétés. Nous pouvons
maintenant énoncer une version très générale du premier théorème
pour les théories du langage arithmétique~$\Lang L_A$.

\begin{thm}
Si $\Gamma$ est une théorie cohérente et !!{axiomatizable}
dans~$\Lang L_A$ qui !!{represents} toutes les fonctions calculables
et toutes les relations !!{decidable}s, alors $\Gamma$ n'est pas
complète.
\end{thm}

Dire que $\Gamma$ n'est pas complète signifie qu'il existe au moins
un !!{sentence}~$!A$ tel que $\Gamma\Proves/ !A$ et
$\Gamma\Proves/\lnot !A$. Un tel !!{sentence} est dit
\emph{indépendant} de~$\Gamma$. Nous pourrons montrer assez vite
qu'il en existe. La force de la méthode de Gödel tient toutefois à
ce qu'elle construit explicitement un !!{sentence}~$!G_\Gamma$,
appelé \emph{énoncé de Gödel} pour~$\Gamma$, qui n'est pas
démontrable : dans~$\Gamma$, $!G_\Gamma$ équivaut à l'affirmation
que $!G_\Gamma$ n'est pas démontrable. Pour construire un énoncé indépendant
sous la seule hypothèse de cohérence, on utilise la variante
de Rosser, qui construit un autre énoncé ; la non-réfutabilité de l'énoncé original de Gödel demande une hypothèse
plus forte.\footnote{Précision éditoriale : la source passe sans qualification de l'énoncé de Gödel à l'indépendance. Nous distinguons les deux constructions et leurs hypothèses.} La construction de l'énoncé de Gödel est \emph{effective} : à partir d'une
axiomatisation de~$\Gamma$ et d'une description du système de
!!{derivation}, la preuve fournit une méthode pour écrire
effectivement~$!G_\Gamma$.

La construction exige que le langage~$\Lang L_A$ puisse exprimer
des propriétés des termes et des !!{formula}s de~$\Lang L_A$
lui-même, ainsi que des opérations sur ces objets. Il faut par
exemple pouvoir exprimer « $!A$ est un !!{sentence} »,
« $\delta$ est une !!{derivation} de~$!A$ », ou l'opération
$\Subst{!A}{t}{x}$. Pour cela, il faut, premièrement, coder les
symboles et leurs suites finies, dont sont formés termes,
!!{formula}s et !!{derivation}s, par des nombres naturels : ce sont
les objets dont $\Lang L_A$ peut parler. Deuxièmement, ce codage
doit permettre à~$\Gamma$ de démontrer les faits pertinents. Nous
devrons donc vérifier que les propriétés en jeu correspondent à des
propriétés !!{decidable}s des nombres naturels et que les opérations
correspondent à des fonctions calculables. C'est ce qu'on appelle
l'\emph{arithmétisation de la syntaxe}.

Il faut aussi examiner la condition selon
laquelle $\Gamma$ est « assez forte », c'est-à-dire !!{represents}
toutes les fonctions calculables et les relations !!{decidable}s.
Elle suppose une définition précise de « calculable ». On peut en
donner plusieurs, par exemple au moyen des machines de Turing ou des
programmes écrits dans un langage de programmation général. Mais,
puisque nous voulons représenter ces fonctions et relations dans
une théorie du langage~$\Lang L_A$, il est préférable de choisir
une définition simple de la calculabilité des fonctions numériques :
celle des \emph{fonctions récursives}, déjà étudiées dans la partie sur la calculabilité.
Après l'arithmétisation de la syntaxe, nous montrerons que $\Th{Q}$ représente déjà toutes les
fonctions et relations récursives.\footnote{Adaptation éditoriale : les fonctions récursives précèdent cette partie dans le lecteur ; les chapitres qui suivent traitent d'abord de l'arithmétisation, puis de la représentation dans~$\Th{Q}$.} Nous pourrons alors appliquer le
théorème d'incomplétude à des théories précises comme $\Th{Q}$ et
$\Th{PA}$, puisqu'elles satisfont la condition de force requise.

L'arithmétisation de la syntaxe aboutira à une !!{formula}
$\OProv[\Gamma](x)$ qui, grâce au codage numérique des !!{formula}s,
exprime la démontrabilité à partir des axiomes de~$\Gamma$. Si
$!A$ a pour code le nombre~$n$ et si $\Gamma\Proves !A$, alors
$\Gamma\Proves\Prov[\Gamma](\num{n})$. Ce \emph{prédicat de
démontrabilité} permet aussi, en un certain sens, d'exprimer la
cohérence de~$\Gamma$ par un !!{sentence} de~$\Lang L_A$.
Choisissons pour~$n$ le code d'une contradiction, par exemple
de~$\lfalse$ : l'\emph{énoncé de cohérence} de~$\Gamma$ sera
$\lnot\Prov[\Gamma](\num{n})$. Le second théorème d'incomplétude
affirme qu'une théorie cohérente et !!{axiomatizable} ne démontre pas
non plus son propre énoncé de cohérence, sous des hypothèses un peu
plus exigeantes que la seule représentation de toutes les fonctions
calculables et relations !!{decidable}s, notamment les conditions de démontrabilité du prédicat usuel. Nous montrerons que
$\Th{PA}$ satisfait ces hypothèses.

\end{document}
